\documentclass[11pt,a4paper]{article}
\usepackage[T1]{fontenc}
\usepackage[margin=1in]{geometry}
\usepackage{amsmath,amssymb,amsthm}
\usepackage{booktabs}
\usepackage{xcolor}
\usepackage[colorlinks=true,linkcolor=blue!60!black,citecolor=blue!60!black,
            urlcolor=blue!60!black]{hyperref}
\usepackage{microtype}
\usepackage{enumitem}
\usepackage[most]{tcolorbox}

\newcommand{\lean}[1]{\texttt{\small #1}}
\newtcolorbox{keybox}[1][]{colback=blue!3,colframe=blue!45!black,
  boxrule=0.5pt,left=4pt,right=4pt,top=3pt,bottom=3pt,#1}
\newtcolorbox{failbox}[1][]{colback=red!3,colframe=red!45!black,
  boxrule=0.5pt,left=4pt,right=4pt,top=3pt,bottom=3pt,#1}

\title{\textbf{A Lean~4 tribute to Cédric Villani:}\\
\large a discrete Brunn--Minkowski inequality, fully formalized,\\
and a stability criterion that turned out not to be formalizable yet}

\author{Xavier Callens\\
\small SocrateAI-Scientific-QuantumFluids\\
\small \url{https://github.com/xaviercallens/SocrateAI-Scientific-QuantumFluids}}
\date{September 2026}

\begin{document}
\maketitle

\begin{abstract}
\noindent
We report a small Lean~4 formalization offered as a tribute to Cédric Villani, built from a literature
survey of his freely available 2025--2026 work and lecture notes rather than from the one result of his
that already has a machine-checked proof (nonlinear Landau damping, formalized independently by
Bedrossian). The survey found that Mathlib already proves the Kullback--Leibler Data Processing
Inequality for a Markov kernel, so we state its three-line specialization to an invariant measure --
the discrete, relative-entropy form of Boltzmann's H-theorem -- explicitly \emph{not} as a
contribution. The headline result is a full, \lean{sorry}-free proof of Theorem~1 of
Ollivier and Villani's \emph{curved Brunn--Minkowski inequality on the discrete hypercube}
(SIAM J.\ Discrete Math.\ 2012) at zero curvature: for nonempty $A,B$ in the Hamming cube with midpoint
set $M$, $\#A\cdot\#B\le(\#M)^2$, by the crossover-coding injection $A\times B\hookrightarrow M\times M$
the paper describes. We also report a companion literature gate, run before any code, on a separate proposal from this
library's kinetic-theory side: to identify Pomeranchuk's thermodynamic stability criterion for a Landau
Fermi liquid with a Penrose-type dynamical one for the two-dimensional Landau kinetic equation. The
literature establishes that correspondence for an analogous three-dimensional system, but not, so far as
we found, for the two-dimensional case; and the numerical parameters needed even to state the
two-dimensional version are not published anywhere we could reach. We stopped rather than derive either
from scratch, and report why.
\end{abstract}

\section{Introduction}

This library already leans on three large machine-checked developments: OpenAI's Navier--Stokes/Euler
formalization~\cite{OpenAINSE}, Anthropic's formalization of Fermat's Last Theorem, evaluated as a
source of general analysis~\cite{AnthropicFLT}, and Bedrossian's Lean~4 formalization of nonlinear
Landau damping~\cite{Bedrossian2026}, which subsumes one small lemma we had proved
independently~\cite{CallensKinetic}. Asked to build a foundation formalizing Cédric Villani's own work
as a tribute, the obvious first move was \emph{not} to duplicate that last result, and the second was to
find out what, among his output, is both self-contained and not already done.

\begin{keybox}
\textbf{What is and is not claimed.} \emph{New:} a full proof of Ollivier--Villani's Theorem~1 at zero
curvature (\S\ref{sec:ov}). \emph{Not a contribution, stated as such:} a corollary of Mathlib's own
Data Processing Inequality (\S\ref{sec:kl}). \emph{Not attempted, named:} the curved case of the same
theorem, and everything listed in \S\ref{sec:notdone}. \emph{Stopped at a literature gate, before any
code:} a proposed link between the Pomeranchuk and Penrose stability criteria (\S\ref{sec:gate}).
\end{keybox}

\section{What is already there}
\label{sec:survey}

We read thirteen freely available items: Villani's two solo papers of 2025--2026, \emph{Fisher
Information in Kinetic Theory}~\cite{Villani2025} and \emph{On the kinetic Fokker--Planck equation in
curved geometry}~\cite{Villani2026}; Imbert--Silvestre--Villani~\cite{ISV2024}; Mouhot--Villani on
Landau damping~\cite{MouhotVillani2011} and its short notes; \emph{Hypocoercivity}~\cite{Villani2009};
Desvillettes--Mouhot--Villani on Cercignani's conjecture~\cite{DMV2011}; Lott--Villani on Ricci
curvature via optimal transport~\cite{LottVillani2009,LottVillani2007}; Ollivier--Villani on a discrete
Brunn--Minkowski inequality~\cite{OllivierVillani2012}; \emph{Optimal Transport, Old and New}; and
\emph{H-theorem and beyond}.

Against this, and against Mathlib as pinned by this project (not a different checkout, and not from
memory), we found: \lean{InformationTheory.klDiv} together with the Kullback--Leibler Data Processing
Inequality for a measurable map, a sub-$\sigma$-algebra, and a Markov kernel
(\lean{klDiv\_comp\_right\_le}), Gibbs' inequality (\lean{mul\_log\_le\_klDiv}), and the chain rule; the
L\'evy--Prokhorov metric and Prokhorov's theorem; Gaussian measures and Brownian motion. We did
\emph{not} find Wasserstein distance, Kantorovich--Rubinstein duality, a log-Sobolev or Talagrand
inequality (the name occurs only in a bibliography entry), Bakry--\'Emery theory, Otto calculus, the Boltzmann equation, Fisher information, a Poincar\'e
inequality, a heat kernel, the Fokker--Planck equation, or Ricci curvature via optimal transport.
Outside Mathlib, J.~K.~Miller's Lean~4 development of mean-field Vlasov theory (arXiv:2607.08986) proves
$W_1$, couplings, the gluing lemma and Kantorovich--Rubinstein duality for Polish spaces from Mathlib
alone -- the natural next dependency for anyone formalizing optimal transport in this ecosystem, though
we do not use it here.

\section{A corollary, explicitly not a contribution}
\label{sec:kl}

For a Markov kernel $\kappa$ and a measure $\nu$ invariant under it, $\kappa\circ_m\nu=\nu$, the Data
Processing Inequality applied to $\mu$ and $\nu$ gives
$\mathrm{klDiv}(\kappa\circ_m\mu,\nu)=\mathrm{klDiv}(\kappa\circ_m\mu,\kappa\circ_m\nu)\le
\mathrm{klDiv}(\mu,\nu)$: relative entropy to an invariant measure is nonincreasing along a Markov
chain -- Boltzmann's H-theorem in the form Cover and Thomas give it~\cite{CoverThomas2006}, and the discrete skeleton of the
entropy methods Villani surveys in \emph{H-theorem and beyond}. The Lean proof is three lines,
\lean{klDiv\_nonincreasing\_to\_invariant} in \lean{Villani.lean}; the mathematical content is entirely
Mathlib's own \lean{klDiv\_comp\_right\_le}. We include it because it is the right first step in this
direction and because it connects this file to \lean{ZeroSound.lean}'s Fermi-liquid kinetic theory
inside one project, not because it is ours.

\section{Ollivier--Villani's Theorem~1, at zero curvature}
\label{sec:ov}

Ollivier and Villani compare two notions of discrete Ricci curvature on the hypercube
$X=\{0,1\}^N$ with the Hamming metric~\cite{OllivierVillani2012}. Along the way they prove a discrete
analogue of the classical Brunn--Minkowski inequality: for nonempty $A,B\subseteq X$, with $M$ the set
of \emph{midpoints} of $A$ and $B$ (a point $m$ with $d(m,a)+d(m,b)=d(a,b)$ and $d(m,a)$ within $1$ of
$d(a,b)/2$, for some $a\in A,b\in B$),
\begin{equation}
\ln\#M\ \ge\ \tfrac12\ln\#A+\tfrac12\ln\#B+\tfrac{K}{8}\,d(A,B)^2,\qquad K=\tfrac1{2N}.
\label{eq:t1}
\end{equation}
We formalize the $K=0$ case, $\#A\cdot\#B\le(\#M)^2$, which the paper itself isolates as \S3
(``Brunn--Minkowski inequality without curvature'') before adding the curvature term. Its own proof is an
explicit injection $A\times B\hookrightarrow M\times M$, which we reproduce faithfully -- read from the
paper, not guessed, and re-derived by hand before a line of Lean was written.

\paragraph{The injection.} For $a,b\in X$ let $D(a,b)=\{i:a_i\ne b_i\}$, $r=\#D(a,b)$. For
$c\subseteq D(a,b)$ define the \emph{crossover} $\varphi_c(a,b)_i=a_i$ if $i\in c$, else $b_i$; this is a
midpoint of $a,b$ exactly when $\#c$ is within $1$ of $r/2$. For a \emph{canonical} choice --
$c_r(a,b)$, the $\lfloor r/2\rfloor$ smallest indices of $D(a,b)$ under the order on $\mathrm{Fin}\,N$,
a function of $D(a,b)$ alone -- the pair
$\Phi(a,b)=(\varphi_{c_r}(a,b),\varphi_{D(a,b)\setminus c_r}(a,b))$ is recoverable: one checks directly
that $D(\varphi_c(a,b),\varphi_{D(a,b)\setminus c}(a,b))=D(a,b)$ for \emph{any} $c\subseteq D(a,b)$, so
$r$ is read off the two output midpoints as their own Hamming distance, which pins down $c_r$ again, and
then each coordinate of $a,b$ is read off the two midpoints by whether it lies in $c_r$,
$D(a,b)\setminus c_r$, or neither. This makes $\Phi$ injective, hence
$\#(A\times B)\le\#(M\times M)$.

\begin{table}[h]\centering\small
\begin{tabular}{@{}ll@{}}
\toprule
\lean{crossover}, \lean{diff\_crossover\_left/right} & the crossover map and its two distances \\
\lean{isMidpoint\_crossover} & any size-appropriate crossover is a genuine midpoint \\
\lean{takeHalf}, \lean{card\_takeHalf} & the canonical half, via Mathlib's \lean{Finset.orderEmbOfFin} \\
\lean{encode}, \lean{encode\_fst/snd\_isMidpoint} & $\Phi$, and that both components are midpoints \\
\lean{diff\_crossover\_crossover\_compl} & $D(\Phi(a,b))=D(a,b)$, for any crossover split \\
\lean{encode\_injective} & $\Phi$ is injective (the case split above) \\
\lean{card\_sq\_ge} & \textbf{$\#A\cdot\#B\le(\#M)^2$} \\
\bottomrule
\end{tabular}
\caption{The load-bearing steps from the definitions to Theorem~1 at $K=0$. The module's axiom audit
covers 16 theorems: 15 for this chain (including unfolding and distance lemmas not shown) and the
entropy corollary of \S\ref{sec:kl}. Standard axiom
footprint \{\lean{propext}, \lean{Classical.choice}, \lean{Quot.sound}\}; accepted by both the Lean
kernel and the independent \lean{nanoda} kernel~\cite{Comparator}; three negative controls (a stricter
bound, a relaxed midpoint inequality, a strict Data Processing Inequality) all fail to compile.}
\end{table}

Beyond first principles the proof uses only a set-difference cardinality identity, Mathlib's canonical
order embedding of a finite subset, and the cardinality bound for an injection between finite sets:
nothing in it needed machinery this project did not already have on hand for other reasons.

\section{Not attempted, and why}
\label{sec:notdone}

\begin{itemize}[leftmargin=*,itemsep=2pt]
\item \textbf{The curved case of the same theorem}, \eqref{eq:t1} with $K=1/(2N)$. The paper's own proof
needs concentration of measure in the symmetric group $S_r$ (its Lemma~4, via a quotienting argument
onto the space of ``crossovers'', its Proposition~5) -- machinery we did not build. Left as a named open
target rather than silently dropped.
\item \textbf{The Fisher-information monotonicity results of~\cite{Villani2025,ISV2024}}: Fisher
information is nonincreasing along the spatially homogeneous Boltzmann equation. Tensorisation
($I(f\otimes f)=2I(f)$) and the Cram\'er--Rao inequality look tractable (Mathlib has the needed
integration-by-parts machinery on $\mathbb R^d$); the full monotonicity theorem is a genuine analytic
argument, not attempted.
\item \textbf{A discrete H-theorem for a finite-velocity collision kernel}, by the classical four-fold
symmetrisation of the entropy production, was our first plan, abandoned mid-proof when \S\ref{sec:kl}'s
three-line route through existing Mathlib content turned out to give a sharper and more general fact
directly.
\item \textbf{Everything genuinely out of reach today}: the continuous collisional Boltzmann H-theorem
(no entropy production functional in Mathlib), hypocoercivity in Villani's abstract Hilbert-space
form~\cite{Villani2009}, log-Sobolev and Talagrand inequalities (no Wasserstein distance, no Gaussian
concentration machinery), the Lott--Villani--Sturm curvature-dimension
condition~\cite{LottVillani2009,LottVillani2007}, and nonlinear Landau damping, already
done~\cite{Bedrossian2026}.
\end{itemize}

\section{A literature gate that stopped a proposal before it was code}
\label{sec:gate}

This library's kinetic-theory side had independently proved, in Lean, that undamped zero sound in a
Fermi liquid with a single Landau parameter $F_0^s$ exists if and only if $F_0^s>0$, in both two and
three dimensions~\cite{CallensKinetic}. A natural next step, chosen to bring the Villani side and the
Godfrin side of this project together, was to formalize a stronger, dynamical statement: that the
\emph{linearized} Landau kinetic equation is stable -- no root of its dispersion relation at
$\mathrm{Im}\,\omega>0$ -- if and only if $F_0^s>-1$, the range where Pomeranchuk's thermodynamic
stability criterion holds, and to enumerate the corresponding roots with certified interval arithmetic
for real $^3$He at the pressures Godfrin et al.\ measured, in both bulk and the two-dimensional film
their neutron-scattering experiment used~\cite{Godfrin2012}.

Following this project's rule to search before writing code, we ran that search first.

\begin{failbox}
\textbf{Stopped.} The thermodynamic Pomeranchuk criterion $F_0^s>-1$ is standard and well
sourced~\cite{Pomeranchuk1958}. Its equivalence to the dynamical, Penrose-type statement is not an
established theorem for the case at hand. The closest match,
Kolomeitsev and Voskresensky~\cite{KolomeitsevVoskresensky2016}, obtains exactly that correspondence by
explicit dispersion analysis -- but for a three-dimensional scalar-interaction nuclear-matter model, not
named after Penrose, and not the two-dimensional case this project targets. Formalizing the statement as
we had proposed it would need an \emph{original} two-dimensional dynamical derivation, which is outside
what this project formalizes. Separately, the numerical parameters the proposal needed even to run --
areal-density-dependent Landau parameters for the two-dimensional $^3$He film Godfrin et al.\ actually
measured -- are not available to us: their Nature paper is paywalled, and no companion paper surfaced a
usable table within the search budget. Both halves of the proposal failed their own pre-registered kill
criteria, and neither was implemented.
\end{failbox}

We report this because a stopped proposal, with the reason recorded, is data about the method, not an
absence of one: three prior topological proposals in this project were likewise stopped or refuted after
being pre-registered~\cite{CallensLean}, and the discipline that catches them is the same one that let
\S\ref{sec:ov} be checked with confidence once it was written.

\section{Conclusion}

A literature survey aimed at ``formalize something of Villani's'' found more value in what was already
proved elsewhere (Bedrossian) or already in the library (Mathlib's Data Processing Inequality) than in
inventing a target; the one genuinely new piece of formal content here is a co-authored theorem old
enough, and self-contained enough, that its proof could be read from the paper and checked, rather than
guessed at. The stopped proposal of \S\ref{sec:gate} is offered as the same kind of evidence this
project's other negative results are: not a discovery, but a fact about where the discipline's cost is
paid -- in a search, before any code, rather than in a retraction after.

\paragraph{Availability.} Source, the pre-registration and literature-gate record for
\S\ref{sec:gate}, and every claim's entry in the project ledger:
\url{https://github.com/xaviercallens/SocrateAI-Scientific-QuantumFluids}
(\lean{lean\_src/Villani.lean}, \texttt{docs/designs/ZERO\_SOUND\_LANDAU\_DAMPING\_PROPOSAL.md}).

\bibliographystyle{unsrt}
\bibliography{refs_villani}
\end{document}
